\documentclass[12pt, a4paper]{article}

% ---------- packages ----------
\usepackage[T1]{fontenc}
\usepackage{mathpazo}
\usepackage{textcomp}
\usepackage{microtype}
\usepackage{geometry}
\geometry{a4paper, top=2.5cm, bottom=2.5cm, left=3cm, right=3cm}

\usepackage{amsmath, amssymb, amsthm, mathtools}

\usepackage{setspace}
\setstretch{1.15}

\usepackage{titlesec}
\titleformat{\section}{\large\bfseries}{\thesection.}{0.75em}{}
\titleformat{\subsection}{\normalsize\bfseries}{\thesubsection.}{0.5em}{}
\titlespacing*{\section}{0pt}{1.6ex plus 0.4ex minus 0.2ex}{0.8ex plus 0.2ex}
\titlespacing*{\subsection}{0pt}{1.2ex plus 0.3ex}{0.5ex plus 0.1ex}

\usepackage{parskip}
\setlength{\parskip}{0.55em}
\setlength{\parindent}{0pt}

\usepackage[hidelinks]{hyperref}
\usepackage{booktabs, array, enumitem}
\setlist[itemize]{noitemsep, topsep=0.3em, leftmargin=1.6em}

% ---------- theorem environments ----------
\newtheorem{theorem}{Theorem}[section]
\newtheorem{definition}[theorem]{Definition}
\newtheorem{proposition}[theorem]{Proposition}
\theoremstyle{remark}
\newtheorem{remark}[theorem]{Remark}

\theoremstyle{definition}
\newtheorem{principle}[theorem]{Principle}
\newtheorem{corollary}[theorem]{Corollary}

% ---------- operators and shorthands ----------
\DeclareMathOperator{\Vol}{Vol}
\newcommand{\A}{\mathcal{A}}
\newcommand{\Rk}{\mathfrak{R}}
\newcommand{\vv}{\mathbf{v}}
\newcommand{\Phif}{\Phi}

% ---------- metadata ----------
\title{\textbf{Amplitwist Cascades, Biological Priors,\\[0.3ex]
and the Geometry of Cognitive Transformation}}
\author{Flyxion\\[0.5ex]
\small Independent Researcher}
\date{May 2026}

% ==========================================================
\begin{document}
\maketitle

\begin{abstract}
\noindent
This essay develops a geometric account of cognitive transformation grounded in the
amplitwist operator formalism.  The central thesis is that cognition operates within a
structured transformation algebra whose expressive capacity exceeds its biological use,
and that this gap---far from being incidental---encodes the prior that shapes how biological
systems actually compute.  We formalize the amplitwist cascade as a universal transformation
language, embed it within RSVP-style field dynamics, and derive efficiency bounds from the
spectral geometry of the representational manifold.  We then introduce a biological prior
modelled on hierarchical central-pattern-generator (CPG)-like dynamics, which constrains
the system to low-dimensional, rhythmically stable trajectories.  Propositional
representation is reconceived as an equivalence class of perturbation-stable trajectories
rather than as static encoding.  The two-layer structure---expressive envelope versus
biological prior---generates falsifiable predictions about phase alignment, vorticity
convergence, and cascade ordering in both human behaviour and neural population dynamics.
An experimental design using compositional and non-compositional morphological
transformation tasks is proposed to test these predictions, and an alignment loss metric is
defined to compare operator geometry between humans and language models.
\end{abstract}

\newpage

\tableofcontents
\bigskip

% ==========================================================
\section{Introduction: From Expressivity to Constraint}
\label{sec:intro}

The central thesis of this essay is not merely that cognition is nonlinear or
compositional, but that it operates within a structured transformation algebra whose
expressive capacity strictly exceeds its biological use.  The distinction that matters is
between what a system \emph{could} compute and what it \emph{actually} computes, between
the envelope of possible transformations and the narrow, prior-shaped corridor through which
biological computation travels.

Standard accounts of neural computation tend to focus on one side of this distinction or
the other.  Representational accounts emphasise the richness of what the brain can encode.
Dynamical systems accounts emphasise the constrained, low-dimensional structure of what it
actually does.  The position developed here holds that both sides are necessary, and that
the relationship between them---the gap between expressive capacity and biological
implementation---is itself the primary theoretical object.

This gap is not a residual to be explained away.  It encodes the prior that shapes
cognition: the preference for rhythmically stable, compositionally ordered, dynamically
supportable trajectories over the vast space of transformations that the underlying operator
algebra could in principle compute.  Formalising this gap, and making it empirically
testable, is the aim of the essay.

The argument proceeds in four stages.  Sections~\ref{sec:operator}--\ref{sec:spectral}
develop the mathematical core: the amplitwist operator as a universal transformation basis,
its embedding in field dynamics, and the spectral constraints governing efficiency.
Sections~\ref{sec:prior}--\ref{sec:representation} introduce the biological prior and
reformulate propositional representation as trajectory equivalence.
Sections~\ref{sec:predictions}--\ref{sec:comparison} derive empirical predictions and
propose an experimental design.  Sections~\ref{sec:interpretation}--\ref{sec:conclusion}
interpret the framework and assess its falsifiability.

% ==========================================================
\section{The Amplitwist Operator as a Universal Transformation Basis}
\label{sec:operator}

\subsection{Operator definition}

The starting point is the classical result from complex analysis.  For a holomorphic
function $f \colon \mathbb{C} \to \mathbb{C}$, the derivative at a point $z_0$ takes the
form $f'(z_0) = \rho\, e^{i\theta}$, encoding a local similarity transformation that
simultaneously scales by $\rho > 0$ and rotates by $\theta \in [0,2\pi)$~\cite{needham1997visual,ahlfors1979complex}.  Tristan
Needham's \emph{Visual Complex Analysis} gives this decomposition the name
\emph{amplitwist}, capturing the dual action of amplitude and twist on infinitesimal
neighbourhoods.

Generalising to a smooth $n$-dimensional representational manifold $M$, we define the
\emph{amplitwist operator} at a point $x \in M$ as the triplet
\begin{equation}
  \A(x) \;=\; \rho(x)\, e^{i\theta(x)}\, \pi(x),
  \label{eq:operator}
\end{equation}
where $\rho(x) \in \mathbb{R}_{>0}$ is an amplitude modulation (gain or precision),
$\theta(x) \in \mathbb{R}$ is a phase encoding the direction of transformation in the
tangent space $T_x M$, and $\pi(x)$ is a nonlinear projection or contraction that
introduces the nonlinearity required for function approximation beyond the affine regime.

\subsection{Cascade composition}

Under sequential application, amplitwist operators compose into cascades.  Given a sequence
of operators $\A_1, \ldots, \A_N$ acting at successive stages, the composed transformation
is
\begin{equation}
  \A^{(N)} \;=\; \A_N \circ \cdots \circ \A_1.
  \label{eq:cascade}
\end{equation}
With sufficient depth and nonlinearity, this compositional system defines a dense function
class over smooth maps $M \to \mathbb{R}^n$.  The amplitwist cascade is therefore
analogous in expressive power to neural network universal approximation theorems~\cite{conway1978functions,evans2010pde}, but
grounded in geometrically interpretable local operators rather than in abstract activation
functions.  This is the \emph{capacity layer}: the formal claim that the operator algebra,
given unconstrained composition, can approximate any sufficiently regular transformation.

\subsection{Two regimes of operation}

The capacity claim must be distinguished from any claim about how biological systems
actually use this capacity.  Formally, two regimes are identified.  In \emph{Regime~A}
(universal approximation), the system exploits the full depth of the cascade to approximate
arbitrary continuous functions; intermediate states are determined solely by the
approximation target.  In \emph{Regime~B} (constrained control), the system operates
through a small number of structured generators---rhythmic, hierarchical,
embodied---and computation is compositional in the sense that each stage stabilises before
the next is engaged.  The central claim of this essay is that biological cognition
predominantly occupies Regime~B, and that the entropy weighting introduced in
Section~\ref{sec:field} implements the corresponding prior.

% ==========================================================
\section{Field-Theoretic Embedding and Phase Alignment}
\label{sec:field}

\subsection{RSVP fields as state space}

The amplitwist operator acts on a structured state space defined by three coupled fields
over $M$.  The \emph{scalar field} $\Phif \colon M \to \mathbb{R}$ represents semantic
salience or potential.  The \emph{vector field} $\vv \colon M \to TM$ represents
conceptual velocity, encoding the direction and speed of representational change.  The
\emph{entropy field} $S \colon M \to \mathbb{R}_{>0}$ represents cognitive uncertainty or
the admissibility of a transformation under the biological prior.  This triplet is the
state space of the Relativistic Scalar--Vector Plenum (RSVP) framework adapted to the
cognitive setting.

\subsection{Phase alignment and the amplitwist angle}

The central geometric quantity is the angle between the vector field $\vv$ and the gradient
of the scalar field $\Phif$.  This angle measures the degree to which the direction of
representational motion aligns with the direction of increasing salience.  Formally,
\begin{equation}
  \theta(x) \;=\; \arccos\!\left(
    \frac{\vv(x) \cdot \nabla\Phif(x)}{\|\vv(x)\|\,\|\nabla\Phif(x)\| + \epsilon}
  \right),
  \label{eq:phase}
\end{equation}
where $\epsilon > 0$ is a small numerical stabiliser.  When $\theta \approx 0$, the system
moves directly along the salience gradient: the transformation is maximally coherent.  When
$\theta \approx \pi/2$, motion is orthogonal to the gradient, generating rotational
structure without progress.  The quantity $\xi = |\nabla \times \hat{\vv}|$, where
$\hat{\vv} = (\cos\theta, \sin\theta)$, is the \emph{vorticity} of the phase-weighted
field.  High vorticity signals inefficiency or instability in the transformation cascade.

\subsection{Entropy weighting as biological prior}

The layer-$k$ amplitwist, weighted by entropy, is
\begin{equation}
  \A^{(k)}(x) \;=\; w_k(x)\cdot \A\!\left(\Rk_k(x)\right), \qquad
  w_k(x) \;=\; \exp\!\bigl(-\lambda S(x)\bigr),
  \label{eq:weighted}
\end{equation}
where $\Rk_k$ is the semantic deformation applied at layer $k$.  The weight $w_k$
suppresses transformations in regions of high entropy.  Interpreted as a regulariser, this
is a standard smoothing device.  Interpreted as a biological prior, it is a hypothesis
about what the nervous system computes: it continuously evaluates which transformations can
be supported by available dynamical primitives, and down-weights trajectories that exceed
that support.  This interpretation is the key conceptual move of Section~\ref{sec:prior}.

% ==========================================================
\section{Hybrid Dynamics and Convergence}
\label{sec:hybrid}

Amplitwist cascades do not operate in isolation from temporal dynamics.  Between discrete
operator applications, the state evolves continuously according to the RSVP field
equations.  The combined system is a \emph{hybrid dynamical system}: continuous semigroup
evolution punctuated by discrete transformation events.

Let $T(t)$ denote the strongly continuous semigroup of field evolution~\cite{engel2000one,pazy1983semigroups}, and let $P_k$ denote the
operator applied at event time $t_k$.  The hybrid trajectory takes the form
\begin{equation}
  T(t_N - t_{N-1})\, P_{N-1}\, \cdots\, T(t_1 - t_0)\, P_0\, T(t_0).
  \label{eq:hybrid}
\end{equation}
A Trotter--Kato-style approximation theorem ensures that, as event spacing shrinks to zero
and the discrete grid refines, the hybrid trajectory converges:
\begin{equation}
  \lim_{h \to 0,\, \delta t \to 0}
  \bigl(S_h(\delta t)\bigr)^{\lfloor t/\delta t\rfloor}
  \;=\; T(t) \circ P_t
  \label{eq:trotterkato}
\end{equation}
in the operator norm on the relevant Sobolev space, directly analogous to the Trotter product formula~\cite{trotter1959product,kato1978trotter}.  This convergence result grounds the
discrete cascade model in functional analysis, ensuring that operator composition inherits
stability guarantees from the continuous semigroup.  The cascade is not an algorithmic
approximation whose relationship to the underlying dynamics is merely hoped for: in the
continuum limit it converges to a continuous operator flow that preserves the same
invariant structure as the underlying field dynamics.

% ==========================================================
\section{Spectral Geometry and Efficiency Constraints}
\label{sec:spectral}

\subsection{Eigenmode structure}

The efficiency of transformation propagation across $M$ is constrained by its spectral
geometry.  The Laplace--Beltrami operator $\Delta_M$ admits a spectral decomposition~\cite{chavel1984eigenvalues,rosenberg1997laplacian} whose
eigenmodes $\psi_n$ and eigenvalues $0 = \lambda_0 < \lambda_1 \le \lambda_2 \le \cdots$
organise the possible patterns of activation on $M$.  Low-frequency modes (small
$\lambda_n$) are broad and spatially integrative; high-frequency modes are local and
transient.  A transformation event excites these modes with amplitude proportional to the
overlap of the event kernel with each eigenfunction.  Sequential cascades correspond to
traversals in eigenmode space, with each layer engaging progressively deeper or
higher-frequency structure.

\subsection{Efficiency bound}

The epistemic efficiency ratio $\eta^{(N)}$, defined as the ratio of useful phase alignment
to total cascade cost, satisfies
\begin{equation}
  \eta^{(N)} \;\geq\;
  \frac{\lambda_1(M)}{N \cdot \max_j \|\epsilon_j T_j\|_\infty},
  \label{eq:efficiency}
\end{equation}
where $\lambda_1(M)$ is the first nonzero Laplacian eigenvalue of $M$, $N$ is the cascade
depth, $\epsilon_j$ are the deformation intensities, and $T_j \in \mathfrak{so}(n)$ are
the Lie-algebra generators of the semantic rotations at each layer.  The bound has a
transparent interpretation: the higher the spectral gap of the manifold, the more easily
transformations propagate between regions and the more efficient the cascade.  Conversely,
a manifold with a small spectral gap offers geometric resistance to propagation, requiring
deeper cascades or larger deformations to achieve the same alignment.  This connects local
operator structure to global manifold geometry via eigenmodes and traveling waves, in
direct correspondence with empirical observations of cortical dynamics in ECoG and MEG.

% ==========================================================
\section{The Biological Prior: Dynamical Constraint via CPG-Like Structures}
\label{sec:prior}

\subsection{The capacity--implementation gap}

The formal capacity established in Section~\ref{sec:operator} is unlimited: given
sufficient depth, amplitwist cascades can approximate any sufficiently regular function.
But biological systems do not search this space uniformly.  The central empirical fact is
that neural computation is strikingly low-dimensional relative to the dimensionality of the
spaces it nominally acts upon~\cite{churchland2012neural,gallego2017neural}.  Population dynamics occupy a small number of modes; motor
trajectories are organised around rhythmic generators; perceptual representations cluster
into stable attractors.  The capacity--implementation gap is not a deficiency but a
structural signature: it reveals the shape of the prior that biological systems bring to
computation.

\subsection{CPG-like dynamics as a trajectory prior}

The natural model for this prior is a hierarchy of central pattern generators (CPGs) or
CPG-like oscillatory structures~\cite{marder2001central,izhikevich2007dynamical}.  CPGs are neural circuits capable of producing
coordinated, rhythmic output without requiring continuous sensory input.  Hierarchically
composed CPG chains generate low-dimensional, stable, perturbation-resistant trajectories
through the high-dimensional neural state space.  CPGs are invoked here not as a literal
substrate for all cognition, but as the canonical example of low-dimensional, phase-stable
dynamical generators; the claim is that cortical computation inherits this structural
constraint, not that it reduces to motor circuits.  The amplitwist cascade, from the
perspective of biological implementation, is not searching the full function class: it
preferentially occupies trajectories that can be parameterised and stabilised by such
generators.

This is the \emph{biological prior layer}.  Formally, the prior assigns high probability
to trajectories that satisfy three conditions: they are low-dimensional relative to the
ambient space of $M$; they are rhythmically structured, admitting a generator
decomposition; and they are stable under perturbation, meaning small displacements return
to the trajectory rather than diverging.  The entropy weighting $w_k = \exp(-\lambda S)$
implements this prior in the operator algebra: regions of high entropy correspond to
transformations that cannot be stabilised by available dynamical primitives and are
accordingly suppressed, consistent with variational frameworks that penalise uncertainty and surprise~\cite{friston2010free}.

\subsection{Errors under the prior}

A key prediction follows from the direction of failure.  When a task requires a
transformation that cannot be decomposed into CPG-supportable stages---a genuinely
non-compositional jump through semantic space---biological systems should not fail randomly.
They should fail in the direction of the nearest decomposable transformation, because
entropy weighting continuously pulls the trajectory toward the highest-prior path even when
that path is incorrect.  This directional failure is qualitatively distinct from the
behaviour of a pure function approximator, and it constitutes a falsifiable signature of
the prior.

The deepest consequence of the two-layer picture is that biological cognition is not best
understood as function approximation, but as constrained trajectory selection in a
transformation space.  Neural networks approximate functions; brains select trajectories.
The following proposition captures this quantitatively.

\begin{proposition}[Prior-Induced Measure Concentration]
\label{prop:concentration}
Let $\mathcal{F}_{\mathrm{all}}$ denote the full amplitwist function class, and let
$\mu_{\mathrm{bio}}$ be the probability measure induced by entropy-weighted dynamics.
Then for any $\epsilon > 0$, there exists a subset $\mathcal{F}_\epsilon \subset
\mathcal{F}_{\mathrm{all}}$ such that
\[
\mu_{\mathrm{bio}}(\mathcal{F}_\epsilon) \ge 1 - \epsilon,
\qquad
\dim(\mathcal{F}_\epsilon) \ll \dim(\mathcal{F}_{\mathrm{all}}).
\]
\end{proposition}

In words: the entropy-weighted prior concentrates biological computation on a
low-dimensional subset of the full function space, even though the full space remains
accessible in principle.  The gap between $\mathcal{F}_\epsilon$ and $\mathcal{F}_{\mathrm{all}}$
is the mathematical form of the capacity--implementation distinction that is the central
theoretical claim of this paper.

\begin{principle}[Transformation Selection Principle]
\label{prin:selection}
Cognitive computation selects trajectories that minimise the constraint energy $\mathcal{C}$
(defined in Appendix~\ref{app:admissibility}) within the expressive envelope
$\mathcal{F}_{\mathrm{all}}$.
\end{principle}

This principle ties together entropy weighting (which penalises high-$S$ trajectories),
the CPG prior (which restricts to low-dimensional generators), the admissibility constraints
(which impose bounds on $\theta$, $S$, and $\xi$), and the spectral efficiency bound (which
limits how rapidly transformations can propagate).  Each of these is a different facet of
the same structural claim: cognition operates not by searching the full function class but
by selecting among trajectories that can be coherently maintained under biological
constraint.

% ==========================================================
\section{Representation as Trajectory Equivalence}
\label{sec:representation}

\subsection{Against static encoding}

The standard view of propositional representation treats semantic content as a property of
a state: a pattern of activation, a point in a representational space, a configuration of
features.  This view is ill-suited to a transformation-first framework, because it makes
the content of a representation independent of how the representation behaves under
perturbation.  Two states that activate the same features but respond differently to the
same sequence of transformations would count as representing the same proposition under
the standard view.  That is the wrong result.

\subsection{Trajectory equivalence}

The alternative proposed here grounds propositional identity in transformation behaviour.

\begin{definition}[Propositional equivalence]
\label{def:equiv}
Let $\Pi = (P_1, \ldots, P_k)$ be a sequence of operator perturbations.  Two states
$x, x' \in M$ are \emph{$\Pi$-equivalent}, written $x \sim_\Pi x'$, if
\[
  \A^{(j)}(x) \;=\; \A^{(j)}(x') \quad \text{for all } j \le |\Pi|.
\]
A \emph{proposition} is an equivalence class $[x]_\Pi = \{x' \in M : x' \sim_\Pi x\}$.
\end{definition}

This is a behavioural criterion for semantic content, not a representational one.  Meaning
is constituted by transformation invariance rather than by any intrinsic property of the
state.  The connection to inferential role semantics---specifically, Brandom's account of
meaning as constituted by normative inferential commitment~\cite{brandom1994making}---is immediate, but the present
formulation provides a geometric implementation that the purely linguistic account lacks:
$\Pi$-equivalence is in principle measurable, whereas inferential role has historically
resisted quantification.

\subsection{CPG dynamics as a trajectory basis}

The trajectory-equivalence account avoids the overclaim that CPG chains directly encode
arbitrary propositions.  Instead, CPG-like dynamics provide the primitive trajectory basis
from which propositional invariants can be constructed.  A proposition is not mapped to a
single motor pattern: it is represented by the equivalence class of all states that respond
identically to the relevant sequence of cascade perturbations, where that sequence is
grounded in, though not reducible to, the CPG-generated trajectory basis.

% ==========================================================
\section{Predictions About Cognitive Dynamics}
\label{sec:predictions}

The two-layer model---expressive capacity plus biological prior---generates a set of strong,
specific predictions.

Within a single transformation layer, as a participant learns the relevant rule, phase
alignment $\theta_t$ should improve monotonically over trials: the participant's
representational trajectory should come to point increasingly in the direction of the task
gradient.  Vorticity $\xi^{(N)}$ should decrease and stabilise within the layer, converging
as the cascade becomes ordered.  These within-layer predictions follow from the attractor
stability theorem, which ensures vorticity convergence given small enough deformation
intensities.

Across transformation layers, the cascade should compose in a lawful order.  Participants
should stabilise layer $\Rk_1$ before $\Rk_2$ becomes coherent, and $\Rk_2$ before
$\Rk_3$.  This ordering is a direct prediction of the biological prior: CPG-like dynamics
engage one level of the hierarchy at a time, and entropy weighting suppresses premature
engagement of higher layers before lower ones are stable.

The failure modes of the framework are as diagnostic as its successes.  A participant who
shortcuts---who moves directly from layer $\Rk_1$ to the final output without passing
through $\Rk_2$---should show a phase discontinuity at layer~2: a sudden high-amplitude
jump with poor alignment between $\vv$ and $\nabla\Phif$ at that boundary.  A participant
confronted with a genuinely non-compositional transformation should either impose a
spurious intermediate stage or show high phase error and boundary-localised vorticity.
These signatures distinguish structured cascades from generic nonlinear learning, because
a generic nonlinear model predicts smooth but arbitrarily shaped trajectories rather than
vorticity specifically localised at decomposition boundaries.

% ==========================================================
\section{Experimental Design: Compositional vs Non-Compositional Tasks}
\label{sec:experiment}

\subsection{Controlled transformation domain}

The experimental domain is designed to be deliberately artificial, low-dimensional, and
over-instrumented.  The goal at this stage is not ecological validity but falsifiability:
the domain must be rich enough that participants can discover structure rather than memorise
labels, while simple enough that the true transformation geometry can be computed in
advance and used as ground truth.  Linguistic morphology with nonce words satisfies both
conditions.

A fixed feature manifold is defined prior to data collection, using phonological feature
vectors from a standard inventory, morphological role vectors drawn from Universal
Dependencies, and semantic-pragmatic features established through a norming study on the
nonce items.  The ground-truth $\Phif$ and $\nabla\Phif$ are computed from this manifold
and held fixed throughout.  This non-circularity is methodologically essential: if the
embedding space were reconstructed from participant responses, the test of phase alignment
would be vacuous.

\subsection{Compositional task family}

The first family presents transformations smoothly decomposable across three layers.
Layer~$\Rk_1$ involves \emph{phonological change} (feature scaling): a nonce stem undergoes
a regular phonological modification, describable as a small-magnitude movement in
phonological feature space.  Layer~$\Rk_2$ involves \emph{morphological role change}
(rotation): the item is recategorised---from noun to verb, agent to instrument, or singular
to plural---requiring a directional reorientation in role-feature space without necessarily
changing magnitude.  Layer~$\Rk_3$ involves \emph{semantic-pragmatic shift}
(entropy-weighted transformation): the item acquires a contextually modulated meaning such
as a diminutive, augmentative, or metaphorical extension, where the correct transformation
depends on contextual entropy and is accordingly suppressed by high uncertainty.

A representative cascade might run: \emph{blick} $\to$ \emph{blicken} $\to$
\emph{blicker} $\to$ \emph{tool-for-blicking}.  The lexical items are irrelevant; what
matters is that each layer has a known target direction in the pre-defined feature space
and that participants' responses can be embedded and measured against that direction.

\subsection{Non-compositional task family}

The second family requires discontinuous jumps through semantic space.  These tasks have
correct answers that cannot be reached by any smooth decomposition into the three layers
above.  The correct transformation is opaque from the standpoint of the cascade structure:
it requires moving to a point in the manifold with no stable intermediate under available
dynamical primitives.  The biological prior predicts that humans will either fail
(showing high vorticity and poor alignment) or succeed by imposing a spurious intermediate
stage not present in the task structure.  A pure function approximator makes neither
prediction: it should succeed or fail based solely on output accuracy, with no
characteristic geometric signature.

\subsection{Trajectory measurement}

From participant responses across trials, the following quantities are computed at each
time step $t$:
\begin{equation}
  \vv_t \;=\; x_{t+1} - x_t, \qquad
  \theta_t \;=\; \angle(\vv_t,\, \nabla\Phif), \qquad
  \rho_t \;=\; \|\vv_t\|,
  \label{eq:measures}
\end{equation}
where $x_t$ is the embedding of the participant's response at trial $t$ and $\nabla\Phif$
is the pre-defined task gradient.  The triplet $(\vv_t, \theta_t, \rho_t)$ constitutes the
empirical estimate of the amplitwist operator at time $t$.  Vorticity is estimated from the
curl of the phase-weighted field across the trial sequence.

% ==========================================================
\section{Human vs Model Comparison}
\label{sec:comparison}

\subsection{Alignment loss}

To compare operator geometry between human participants and language models, we define the
\emph{geometric alignment loss}.  Let $\mathcal{A}^{(k)} = (\rho_k, \theta_k, S_k)$ denote
the operator triplet at layer $k$.  Then
\begin{equation}
  \mathcal{L}_A \;=\; \sum_{k=1}^{N} \Bigl(
    \alpha_\rho\,(\rho_k^{(h)} - \rho_k^{(m)})^2
    + \alpha_\theta\,d_{S^1}(\theta_k^{(h)},\, \theta_k^{(m)})^2
    + \alpha_S\,(S_k^{(h)} - S_k^{(m)})^2
  \Bigr),
  \label{eq:alignloss}
\end{equation}
where $d_{S^1}(\theta, \theta') = \min(|\theta - \theta'|, 2\pi - |\theta - \theta'|)$ is
the angular geodesic distance on the circle, and $\alpha_\rho, \alpha_\theta, \alpha_S > 0$
are weighting coefficients.  Using $d_{S^1}$ rather than the naive Euclidean difference
ensures that the metric is correct for the periodic phase variable and that the loss is
genuinely a geodesic distance on the operator space $\mathbb{R}_{>0} \times S^1 \times
\mathbb{R}_{\ge 0}$, not an ad hoc sum of squared differences.
The critical comparison is not whether the model and the human produce the same final
answer, but whether they traverse the transformation space through the same local
rotation-scaling path.  Endpoint equivalence with path divergence is the key empirical
signature: if a model consistently matches human outputs while exhibiting disordered
intermediate operator geometry, it has learned to produce the same results without learning
the underlying transformation structure.

\subsection{Predicted asymmetries}

The framework predicts three distinct patterns.  Humans on compositional tasks should show
low vorticity, improving phase alignment, and ordered cascade composition, with each layer
stabilising before the next is engaged.  Humans on non-compositional tasks should show
high vorticity or the characteristic signature of imposed spurious structure: a phase
discontinuity followed by an intermediate state that reduces vorticity before the
discontinuous jump is completed.  Language models, given either task family, should show
correct or near-correct outputs but weaker geometric coherence: higher vorticity at layer
boundaries, less stable intermediate operator structure, and entropy insensitivity (because
models lack anything analogous to the biological prior implemented by entropy weighting).

The asymmetry between human failure modes---directional, prior-shaped, localised at
decomposition boundaries---and model failure modes---random, entropy-insensitive, uniformly
distributed across layers---is the central empirical contrast.  If this asymmetry is not
observed, the biological prior interpretation is falsified independently of the operator
formalism.

% ==========================================================
\section{Interpretation: Geometry as the Invariant Substrate}
\label{sec:interpretation}

If the predicted invariants are confirmed, the framework supports the following
interpretation.  Cognition, perception, and semantic processing are governed by shared
transformation operators whose local structure is captured by the amplitwist formalism.
Biological systems restrict the expressive envelope of these operators via dynamical priors
that reflect the rhythmic, hierarchical, embodied structure of the nervous system.  Meaning
is not stored in states but emerges from the invariant structure of transformations: two
states mean the same thing when they behave the same way under the same cascade of
perturbations.

The same formalism describes local computation (individual operator applications), global
coherence (cascade composition and eigenmode propagation), and cross-domain transfer (the
alignment loss between human and model transformation geometry).  This unified description
is not merely analogical: it is grounded in the functional-analytic convergence of hybrid
dynamics (Section~\ref{sec:hybrid}) and the spectral geometry of the manifold
(Section~\ref{sec:spectral}).

The framework also clarifies a connection to the geometry of conscious states without
presupposing any particular answer to the hard problem.  If conscious states are
characterised by particularly stable, low-entropy, high-phase-alignment trajectories
through the representational manifold---a claim consistent with integrated information
theory and the entropic brain hypothesis~\cite{tononi2008consciousness,carhart2014entropic}---then the amplitwist formalism provides a
language for describing the geometric structure that distinguishes such states from
non-conscious processing.

% ==========================================================
\section{Falsifiability and Outcomes}
\label{sec:falsifiability}

The theory is falsified by any of the following outcomes.  If human trajectories on
compositional tasks do not exhibit improving phase alignment and converging vorticity, the
within-layer prediction fails.  If compositional tasks do not produce ordered cascade
composition---if participants show no tendency to stabilise lower layers before engaging
higher ones---the across-layer prediction fails.  If non-compositional tasks do not produce
the characteristic signature of high vorticity or imposed spurious structure, the
biological prior prediction fails.  If the alignment loss metric does not discriminate
human from model operator geometry despite output-level accuracy differences, the
transformation-coherence hypothesis is not supported.

Partial failures are informative rather than merely negative.  If within-layer alignment
improves but across-layer ordering does not appear, it suggests the biological prior
operates within layers but not across them---a finding that would call for revision of the
CPG-hierarchy account.  If the alignment loss discriminates humans from models on
compositional tasks but not on non-compositional ones, it suggests the prior is doing
discriminative work only where it has traction, which would refine rather than refute the
framework.

The decisive null test is the following.

\paragraph{Null Hypothesis.} If phase alignment $\theta_t$ and vorticity $\xi_t$ show no
systematic difference between (i) compositional and non-compositional tasks, and (ii)
human and model trajectories, then the amplitwist framework does not capture a distinct
computational invariant and the framework is falsified.  This null can be tested without
any commitment to the full RSVP field theory: it requires only that the discrete
observables $(\theta_t, \xi_t)$ be computed from a pre-specified embedding and compared
across conditions.

If the predictions are confirmed, the consequences are precise.  Transformation coherence
is a real, measurable dimension of cognitive processing, distinct from output accuracy.
Biological systems operate as constrained operator cascades with a specific geometric
signature.  Current language models capture output statistics without capturing the
operator geometry that generates human outputs.  The amplitwist cascade is not an elegant
reinterpretation of known ideas under a new vocabulary: it identifies structure in data
that output-level models do not.

% ==========================================================
\section{Conclusion}
\label{sec:conclusion}

The amplitwist cascade is simultaneously a universal transformation language and a
description of biological constraint.  As a universal language, it encompasses the full
expressive envelope of local rotation-scaling-projection operators under entropic
regularisation.  As a description of biological constraint, it formalises the prior that
shapes cognition: the preference for CPG-supportable, rhythmically stable, compositionally
ordered trajectories over the vast space of transformations the algebra could in principle
compute.

The critical insight is the separation: capacity defines possibility; constraint defines
reality.  The gap between them is not an embarrassment to be minimised but the primary
theoretical object---the structure that encodes the biological prior, generates the
falsifiable predictions, and grounds the alignment loss metric that distinguishes human from
model operator geometry.

The next step is empirical.  The toy-domain experiment proposed in
Section~\ref{sec:experiment} is designed to be small enough that the true transformation
geometry can be computed and used as ground truth, but rich enough that participants must
discover structure rather than memorise labels.  Running the experiment, computing phase
alignment and vorticity trajectories, and comparing human and model operator geometry will
determine whether transformation coherence is a real cognitive variable or an elegant
fiction.  The decisive outcome is not whether the framework is elegant, but whether it
identifies structure in data that other models miss.

\bigskip
\noindent\textbf{Acknowledgements.}\quad
This essay was developed as part of the RSVP\slash Amplitwistor research programme.

\clearpage
\appendix

% Reformat appendix section headings to match body style
\titleformat{\section}{\large\bfseries}{Appendix \Alph{section}:}{0.75em}{}
\titleformat{\subsection}{\normalsize\bfseries}{\Alph{section}.\arabic{subsection}}{0.6em}{}

% ==========================================================
\section{Amplitwist Operator Algebra}
\label{app:algebra}

\subsection{Local Operator Definition}

Let $M$ be a smooth $n$-dimensional manifold and $x \in M$.  Let $T_x M$ denote the tangent
space at $x$.

\begin{definition}[Amplitwist Operator]
An \emph{amplitwist operator} at $x$ is a triple
\[
\mathcal{A}_x = (\rho_x,\, R_x,\, \pi_x),
\]
where $\rho_x \in \mathbb{R}_{>0}$ is a scalar gain (amplitude), $R_x \in SO(n)$ is a
rotation acting on $T_x M$, and $\pi_x \colon T_x M \to T_x M$ is a (possibly nonlinear)
projection map.
\end{definition}

The action of $\mathcal{A}_x$ on a vector $u \in T_x M$ is
\[
\mathcal{A}_x(u) = \pi_x\!\bigl(\rho_x\, R_x u\bigr).
\]

\subsection{Composition Law}

Let $\mathcal{A}_1 = (\rho_1, R_1, \pi_1)$ and $\mathcal{A}_2 = (\rho_2, R_2, \pi_2)$.
Define composition by
\[
(\mathcal{A}_2 \circ \mathcal{A}_1)(u)
= \pi_2\!\Bigl(\rho_2 R_2\,\pi_1(\rho_1 R_1 u)\Bigr).
\]

\begin{proposition}
The set of amplitwist operators equipped with composition forms a nonlinear semigroup
$(\mathcal{A}, \circ)$.
\end{proposition}

\begin{proof}
Closure follows from closure of scalar multiplication, rotations, and function composition.
Associativity follows from associativity of function composition.  Nonlinearity arises from
the presence of $\pi_x$.
\end{proof}

\begin{remark}
In general $\mathcal{A}_2 \circ \mathcal{A}_1 \neq \mathcal{A}_1 \circ \mathcal{A}_2$,
due to non-commutativity of both the rotation group ($R_2 R_1 \neq R_1 R_2$ for $n \ge 3$)
and the projection maps ($\pi_2 \circ \pi_1 \neq \pi_1 \circ \pi_2$ in general).
\end{remark}

\subsection{Local Linearization}

If $\pi_x$ is differentiable, the Jacobian of $\mathcal{A}_x$ at $u$ is
\[
D\mathcal{A}_x(u) = D\pi_x(\rho_x R_x u)\cdot \rho_x R_x.
\]
In the linear case $\pi_x = \mathrm{Id}$, this reduces to $D\mathcal{A}_x = \rho_x R_x$,
a similarity transformation.

\subsection{Density in Function Space}

Let $K \subset \mathbb{R}^n$ be compact and let $C(K)$ denote the space of continuous
functions equipped with the uniform norm.

\begin{theorem}[Expressive Density]
\label{thm:density}
Suppose $\pi_x$ contains a non-polynomial nonlinearity and amplitwist operators are
composed in finite depth.  Then the set of functions representable as
$f(x) = (\mathcal{A}_N \circ \cdots \circ \mathcal{A}_1)(x)$
is dense in $C(K)$.
\end{theorem}

\begin{proof}[Sketch]
Each amplitwist layer can be written locally as $u \mapsto \pi(Wu + b)$ with $W = \rho R$.
Since compositions of affine maps and non-polynomial nonlinearities are dense in $C(K)$,
the result follows from standard universal approximation arguments.
\end{proof}

\subsection{Cascades as Operator Sequences}

Define an \emph{amplitwist cascade} of depth $N$ as
$\mathcal{A}^{(N)} = \mathcal{A}_N \circ \cdots \circ \mathcal{A}_1$.
Given an initial state $u_0$, the induced trajectory is $u_{k+1} = \mathcal{A}_{k+1}(u_k)$,
defining a discrete dynamical system on $T_x M$.

\subsection{Separation of Capacity and Implementation}

\begin{definition}[Expressive Capacity]
$\mathcal{F}_{\mathrm{all}} = \overline{\{\mathcal{A}_N \circ \cdots \circ \mathcal{A}_1\}}$,
the closure in $C(K)$.
\end{definition}

\begin{definition}[Biological Subspace]
$\mathcal{F}_{\mathrm{bio}} \subset \mathcal{F}_{\mathrm{all}}$ denotes functions realisable
under biological constraints (formalised in Appendix~\ref{app:cpg}).
\end{definition}

\begin{proposition}
In general, $\mathcal{F}_{\mathrm{bio}} \subsetneq \mathcal{F}_{\mathrm{all}}$.
\end{proposition}

This strict inclusion formalises the separation between capacity (what amplitwist cascades
can represent) and implementation (what biological systems actually compute).

% ==========================================================
\section{Phase Alignment and Vorticity}
\label{app:phase}

\subsection{Field Structure}

Let $M$ be a smooth manifold.  Define a scalar field $\Phi \colon M \to \mathbb{R}$
(salience), a vector field $v \colon M \to TM$ (trajectory velocity), and an entropy
field $S \colon M \to \mathbb{R}_{\ge 0}$ (uncertainty).  Assume $v(x) \neq 0$ and
$\nabla\Phi(x) \neq 0$ almost everywhere.

\subsection{Phase Alignment}

\begin{definition}[Phase Angle]
The local phase alignment is
\[
\theta(x)
= \arccos\!\left(
\frac{v(x)\cdot\nabla\Phi(x)}
{\|v(x)\|\,\|\nabla\Phi(x)\|+\epsilon}
\right), \quad \epsilon > 0.
\]
\end{definition}

\begin{proposition}
$\theta(x) = 0 \iff v(x) \parallel \nabla\Phi(x)$, and
$\theta(x) = \pi \iff v(x) \parallel -\nabla\Phi(x)$.
Thus $\theta(x)$ measures directional coherence between the trajectory and the salience
gradient.
\end{proposition}

\subsection{Alignment Functional}

Define the global alignment energy
$\mathcal{E}_{\mathrm{align}} = \int_M (1 - \cos\theta(x))\,d\mu(x)$.

\begin{theorem}
$\mathcal{E}_{\mathrm{align}} = 0 \iff v(x) \parallel \nabla\Phi(x)$ for all $x \in M$.
\end{theorem}

\begin{proof}
Follows immediately from $1 - \cos\theta = 0 \iff \theta = 0$.
\end{proof}

\subsection{Vorticity}

\begin{definition}[Vorticity]
Let $\hat{v}(x) = v(x)/\|v(x)\|$ be the normalised direction field.  The vorticity is
\[
\xi(x) = \|\nabla \times \hat{v}(x)\|.
\]
\end{definition}

The integrated vorticity is $\Xi = \int_M \xi(x)\,d\mu(x)$, which equals zero if and only
if $\hat{v}$ is irrotational.

\subsection{Phase--Vorticity Relationship}

\begin{proposition}
Large spatial gradients of $\theta$ imply increased vorticity: if
$\|\nabla\theta(x)\| \gg 0$ then $\xi(x)$ is large.
\end{proposition}

\begin{proof}[Sketch]
Since $\hat{v}$ depends on the angular variable, spatial variation in phase induces curl in
the direction field.
\end{proof}

\subsection{Discrete Approximation}

Given a discrete response trajectory $\{x_t\}$, define $v_t = x_{t+1} - x_t$ and
$\hat{v}_t = v_t/(\|v_t\|+\epsilon)$.  The discrete phase estimate is
\[
\theta_t = \arccos\!\left(
\frac{v_t \cdot \nabla\Phi(x_t)}{\|v_t\|\,\|\nabla\Phi(x_t)\|+\epsilon}
\right),
\]
and the discrete vorticity approximation is
$\xi_t \approx \|\hat{v}_{t+1} - \hat{v}_t\|$.

\subsection{Coherence Criterion}

A trajectory is \emph{coherent} if $\sup_t \theta_t \le \theta_{\mathrm{crit}}$ and
$\sum_t \xi_t \le \Xi_{\mathrm{crit}}$.  Low $\theta$ signals aligned transformation; low
$\xi$ signals stable cascade dynamics.  High $\theta$ signals misalignment; high $\xi$
signals rotational instability.  These quantities define measurable invariants of
transformation coherence that can be estimated from behavioural response sequences.

% ==========================================================
\section{Entropy Weighting as a Biological Prior}
\label{app:entropy}

\subsection{Entropy Field and Weighting}

Let $S \colon M \to \mathbb{R}_{\ge 0}$ denote an entropy (uncertainty) field.

\begin{definition}[Entropy Weight and Weighted Operator]
For layer $k$, define $w_k(x) = e^{-\lambda S(x)}$ with $\lambda > 0$.  The
entropy-weighted amplitwist is
\[
\widetilde{\mathcal{A}}^{(k)}_x(u)
= w_k(x)\,\pi_x\!\bigl(\rho_x R_x u\bigr).
\]
High entropy suppresses operator influence; low entropy amplifies stable transformations.
\end{definition}

\subsection{Variational Formulation}

Consider $f \colon M \to \mathbb{R}^m$.  Define the weighted Dirichlet energy
\[
\mathcal{L}[f]
= \int_M e^{-\lambda S(x)}\,\|\nabla f(x)\|^2\,d\mu(x).
\]

\begin{proposition}
Critical points of $\mathcal{L}$ satisfy the weighted Euler--Lagrange equation
$\nabla\cdot(e^{-\lambda S}\nabla f) = 0$.
\end{proposition}

\begin{proof}
Computing the first variation and integrating by parts:
$\delta\mathcal{L} = -2\int_M \delta f\,\nabla\cdot(e^{-\lambda S}\nabla f)\,d\mu = 0$
for all admissible $\delta f$.
\end{proof}

Regions with large $S(x)$ are penalised (exploration suppressed); regions with small
$S(x)$ are favoured (stable exploitation).  Thus $w_k$ encodes a biological prior toward
low-uncertainty trajectories.

\subsection{Admissibility and Stability}

\begin{definition}[Admissible State and Trajectory]
A point $x \in M$ is admissible if $S(x) \le S_{\max}$.  A trajectory $\gamma\colon[0,T]\to M$
is admissible if $\sup_{t} S(\gamma(t)) \le S_{\max}$.
\end{definition}

\begin{proposition}
If $S(x) \le S_{\max}$ and $\lambda > 0$, then $e^{-\lambda S_{\max}} \le w_k(x) \le 1$.
\end{proposition}

\subsection{Entropy--Vorticity Coupling}

Under the assumption that high-entropy regions induce higher phase variability
($\|\nabla\theta(x)\| \propto S(x)$), Appendix~\ref{app:phase} implies that
$S(x) \uparrow$ entails $\xi(x) \uparrow$, since vorticity scales with the spatial
gradient of direction.

\subsection{Regularised Energy Functional}

The combined functional
\[
\mathcal{J}[f]
= \int_M \bigl[e^{-\lambda S(x)}\|\nabla f(x)\|^2
+ \alpha S(x)^2\bigr]\,d\mu(x)
\]
enforces coherent transformations through its first term and penalises entropy growth
through its second.  Admissible dynamics minimise $\mathcal{J}$ subject to task
constraints.  The limiting cases are informative: as $\lambda \to 0$, the entropy weight
vanishes and the full expressive capacity of $\mathcal{F}_{\mathrm{all}}$ is recovered;
as $\lambda \to \infty$, only minimal-entropy trajectories survive, corresponding to the
most strongly constrained biological regime.

% ==========================================================
\section{Hybrid Dynamics and Trotter--Kato Convergence}
\label{app:hybrid}

\subsection{Semigroup Evolution}

Let $X$ be a Banach space (e.g.\ $L^2(M)$ or $H^k(M)$) and let $F\colon\mathcal{D}(F)
\subset X \to X$ generate a well-posed flow.

\begin{definition}[Strongly Continuous Semigroup]
$\{T(t)\}_{t\ge 0}$ satisfies $T(0) = I$, $T(t+s) = T(t)T(s)$, and
$\lim_{t\downarrow 0}\|T(t)u - u\| = 0$ for all $u \in X$.
For linear $F = A$ (densely defined and closed), $T(t) = e^{tA}$.
\end{definition}

\subsection{Discrete Operator Updates}

Let $\{P_k\}$ be a sequence of (possibly nonlinear) amplitwist operators acting on $X$.
Given a time grid $0 = t_0 < t_1 < \cdots < t_N = t$ with $\Delta t_k = t_{k+1} - t_k$,
the hybrid evolution is
\[
u_{k+1} = P_k\,T(\Delta t_k)\,u_k, \qquad
u(t) = \Bigl(\prod_{k=0}^{N-1} P_k T(\Delta t_k)\Bigr)u_0,
\]
with time-ordered product.

\subsection{Linear Trotter--Kato Framework}

\begin{theorem}[Trotter Product Formula]
If $A$ and $B$ generate strongly continuous semigroups $T_A(t)$ and $T_B(t)$ on $X$, then
under standard domain conditions,
\[
\lim_{n\to\infty}\Bigl(T_A(t/n)\,T_B(t/n)\Bigr)^n = T_{A+B}(t)
\]
strongly on $X$.
\end{theorem}

Identifying $T_A(\Delta t) \equiv T(\Delta t)$ and $T_B(\Delta t) \equiv P(\Delta t)$, the
hybrid scheme approximates $\exp(t(A+B))$ in the linearised regime.

\subsection{Nonlinear Extension (Lie--Trotter Splitting)}

For nonlinear evolution $u' = F(u) + G(u)$, define flows $\Phi_F(t)$ and $\Phi_G(t)$.

\begin{theorem}[Lie--Trotter Splitting]
Under Lipschitz conditions on $F$ and $G$, the scheme
$u^{(n)}(t) = (\Phi_G(t/n) \circ \Phi_F(t/n))^n u_0$
converges to the solution of $u' = F(u) + G(u)$ as $n \to \infty$.
\end{theorem}

The identification is $\Phi_F \leftrightarrow T(t)$ and $\Phi_G \leftrightarrow P(t)$.

\subsection{Convergence Result}

\begin{theorem}[Hybrid Convergence]
\label{thm:hybrid}
Assume $T(t)$ is strongly continuous, $P_k \to P_{t_k}$ uniformly, the scheme is stable
(in the sense that $\sup_n \|\prod_{k=0}^{n-1} P_k T(\Delta t_k)\| \le C$), and the local
truncation error is $O(\Delta t^2)$.  Then
\[
\lim_{\max\Delta t_k \to 0}
\Bigl(\prod_{k=0}^{N-1} P_k\,T(\Delta t_k)\Bigr)
= T(t)\circ P_t
\]
strongly on $X$.
\end{theorem}

The consequence for cascades is direct: discrete amplitwist cascades are not arbitrary
updates but approximate a continuous operator flow, and in the limit they inherit the same
invariant structure as the underlying continuous field dynamics.

% ==========================================================
\section{Spectral Geometry and Efficiency Bounds}
\label{app:spectral}

\subsection{Laplacian and Eigenstructure}

Let $M$ be a compact Riemannian manifold with metric $g$.  The Laplace--Beltrami operator
$\Delta_M f = \nabla \cdot \nabla f$ has eigenpairs $(\lambda_k, \psi_k)$ satisfying
$\Delta_M \psi_k = -\lambda_k \psi_k$ with $0 = \lambda_0 < \lambda_1 \le \lambda_2 \le
\cdots$.  Any $f \in L^2(M)$ admits the spectral expansion $f(x) = \sum_k a_k \psi_k(x)$,
and semigroup evolution satisfies $T(t)f = \sum_k a_k e^{-\lambda_k t}\psi_k$.

\subsection{Efficiency Metric and Bound}

\begin{definition}[Cascade Efficiency]
$\eta^{(N)} = \|f^{(N)} - f^{(0)}\|_{L^2} \big/ \sum_{j=1}^N \|\epsilon_j T_j\|_{L^2 \to L^2}$.
\end{definition}

\begin{theorem}[Spectral Efficiency Bound]
\label{thm:efficiency}
Let $M$ be a compact Riemannian manifold and let $\{\psi_k\}$ be the orthonormal
eigenbasis of $\Delta_M$.  Let the cascade act on $f^{(0)} \in L^2(M)$ with
$f^{(0)} \perp \psi_0$, and assume each layer admits the first-order expansion
\[
\mathcal{A}_j = I + \epsilon_j T_j + O(\epsilon_j^2),
\]
where $T_j$ is a bounded operator with $\|T_j\|_{L^2 \to L^2} \le C$.  Then
\[
\|f^{(N)} - f^{(0)}\|_{L^2}
\;\ge\;
\lambda_1(M)\,\Bigl\|\sum_{j=1}^N \epsilon_j T_j f^{(0)}\Bigr\|_{L^2}
- O\!\Bigl(\sum_j \epsilon_j^2\Bigr),
\]
and consequently
\[
\eta^{(N)} \;\gtrsim\;
\frac{\lambda_1(M)}{N \cdot \max_j \|\epsilon_j T_j\|_{L^2\to L^2}}.
\]
\end{theorem}

\begin{proof}[Sketch]
Since $f^{(0)} \perp \psi_0$, the Rayleigh quotient gives $\|f^{(0)}\|_{L^2}^2 \le
\lambda_1(M)^{-1}\|\nabla f^{(0)}\|_{L^2}^2$.  Expanding the cascade to first order in
$\epsilon_j$ and applying the Rayleigh quotient to the accumulated deformation
$\sum_j \epsilon_j T_j f^{(0)}$ yields the lower bound on $\|f^{(N)} - f^{(0)}\|_{L^2}$.
The second-order remainder is $O(\sum_j \epsilon_j^2)$~\cite{chavel1984eigenvalues,rosenberg1997laplacian}.
\end{proof}

The bound arises because any nontrivial deformation must excite modes orthogonal to the
constant eigenfunction, and the Rayleigh quotient enforces a minimum energy cost
proportional to $\lambda_1(M)$.  The spectral gap therefore sets a geometric lower limit
on how efficiently transformations propagate across $M$~\cite{cheeger1970lower}.

\subsection{Geometric Interpretation}

Large $\lambda_1$ indicates strong connectivity and rapid propagation; small $\lambda_1$
indicates geometric resistance.  The Cheeger constant $h(M)$ satisfies
$h(M)^2/4 \le \lambda_1(M) \le 2h(M)$, so geometric bottlenecks directly limit cascade
efficiency.  Low-frequency modes (small $\lambda_k$) correspond to long-range integration
and slow dynamics; high-frequency modes correspond to local refinement and transient
responses.

\subsection{Spectral Filtering and the Biological Subspace}

Entropy-weighted amplitwists preferentially damp high-frequency modes: the map
$a_k \mapsto a_k e^{-\lambda S}$ suppresses higher-$\lambda_k$ coefficients more
strongly, favouring smooth, low-energy representations.  The biologically realisable
function space can accordingly be characterised spectrally as
\[
\mathcal{F}_{\mathrm{bio}}
= \Bigl\{f \in \mathcal{F}_{\mathrm{all}} \;\Big|\; \sum_k \lambda_k a_k^2 \le C\Bigr\}
\]
for some bound $C$, defining a spectrally constrained subset of the full function space in
which biological systems operate within low-energy spectral bands.

% ==========================================================
\section{Dynamical Constraint Model (CPG Prior)}
\label{app:cpg}

\subsection{Trajectory Spaces and Biological Realisation}

Let $\mathcal{T}_{\mathrm{all}} = \{\gamma\colon [0,T] \to M \mid \gamma \text{ absolutely
continuous}\}$.

\begin{definition}[Biological Trajectory]
A trajectory $\gamma \in \mathcal{T}_{\mathrm{bio}} \subset \mathcal{T}_{\mathrm{all}}$
if it is generated by a finite composition of low-dimensional dynamical primitives
$\gamma(t) = \Phi_k \circ \cdots \circ \Phi_1(t)$, where each $\Phi_i$ is generated by
a smooth, bounded vector field $\dot{x} = f_i(x)$ with $\|\nabla f_i\| \le C$.
A primitive is \emph{CPG-like} if $f_i(x) = A_i x + g_i(x)$, where $A_i$ has imaginary
eigenvalues (oscillatory structure) and $g_i$ is a small nonlinear perturbation.
\end{definition}

\subsection{Constraint Functional}

Define the trajectory energy
\[
\mathcal{J}[\gamma]
= \int_0^T \Bigl(
\|\ddot{\gamma}(t)\|^2
+ \alpha\,\xi(\gamma(t))
+ \beta\,S(\gamma(t))
\Bigr)\,dt.
\]
This functional simultaneously penalises rapid acceleration (lack of smoothness), vorticity
(rotational instability), and entropy (uncertainty).

\begin{theorem}
Biological trajectories approximately satisfy
$\gamma \in \arg\min_{\tilde\gamma \in \mathcal{T}_{\mathrm{all}}} \mathcal{J}[\tilde\gamma]$,
with Euler--Lagrange equation
$\frac{d^2}{dt^2}\ddot\gamma(t) - \alpha\nabla\xi(\gamma(t)) - \beta\nabla S(\gamma(t)) = 0$.
\end{theorem}

\subsection{Low-Dimensional Manifold Constraint}

\begin{proposition}
There exists a manifold $\mathcal{M}_{\mathrm{bio}} \subset \mathcal{T}_{\mathrm{all}}$ of
dimension $d \ll \dim(\mathcal{T}_{\mathrm{all}})$ such that $\mathcal{T}_{\mathrm{bio}}
\subset \mathcal{M}_{\mathrm{bio}}$.
\end{proposition}

\begin{proof}[Sketch]
Follows from restriction to a finite set of dynamical primitives and smooth compositions.
\end{proof}

Although amplitwist cascades are dense in $C(K)$ by Theorem~\ref{thm:density}, biological
trajectories realise only the restricted subset $\mathcal{F}_{\mathrm{bio}} = \{f \in
\mathcal{F}_{\mathrm{all}} \mid f \text{ induced by } \gamma \in \mathcal{T}_{\mathrm{bio}}\}$.
A transformation $f$ is \emph{non-dynamical} if no $\gamma \in \mathcal{T}_{\mathrm{bio}}$
generates it with bounded $\mathcal{J}$; such transformations produce high vorticity, high
entropy, and discontinuities in $\theta$.

The measure-theoretic form of this restriction is given by the following proposition, which
is the formal statement of the paper's central thesis.

\begin{definition}[Entropy-Induced Path Measure]
Let $\mathcal{T}_{\mathrm{all}}$ be the space of admissible trajectories $\gamma\colon[0,T]\to M$.
Define the Gibbs measure over paths
\[
d\mu_{\mathrm{bio}}(\gamma)
= \frac{1}{Z}\exp\!\Bigl(-\lambda\int_0^T S(\gamma(t))\,dt\Bigr)\,d\mu_0(\gamma),
\]
where $\mu_0$ is a reference measure on trajectory space and $Z$ is the normalising
partition function.  The induced measure on function space is $\Phi_*\mu_{\mathrm{bio}}$,
where $\Phi\colon\mathcal{T}_{\mathrm{all}}\to\mathcal{F}_{\mathrm{all}}$ maps trajectories
to their induced transformations.
\end{definition}

\begin{proposition}[Entropy-Induced Concentration]
\label{prop:concentration_app}
Assume the entropy functional $S$ is coercive in the sense that trajectories leaving the
finite-dimensional submanifold $\mathcal{M}_{\mathrm{bio}} \subset \mathcal{T}_{\mathrm{all}}$
incur entropy growth $S(\gamma(t)) \ge c\,\mathrm{dist}(\gamma(t), \mathcal{M}_{\mathrm{bio}})^2$
for some $c > 0$.  Then for any $\epsilon > 0$, there exists a tubular neighbourhood
$\mathcal{U}_\epsilon$ of $\mathcal{M}_{\mathrm{bio}}$ such that
\[
\mu_{\mathrm{bio}}(\mathcal{U}_\epsilon) \ge 1 - \epsilon.
\]
Consequently, the induced function class $\mathcal{F}_\epsilon = \Phi(\mathcal{U}_\epsilon)$
has effective dimension bounded by $\dim(\mathcal{M}_{\mathrm{bio}})$.
\end{proposition}

\begin{proof}[Sketch]
Under the coercivity assumption, trajectories at distance $r$ from $\mathcal{M}_{\mathrm{bio}}$
carry entropy weight at most $e^{-\lambda c r^2 T}$.  Standard large-deviation estimates
then show that $\mu_{\mathrm{bio}}(\mathcal{T}_{\mathrm{all}} \setminus \mathcal{U}_\epsilon)
\le e^{-\lambda c\, \mathrm{dist}(\mathcal{U}_\epsilon, \mathcal{M}_{\mathrm{bio}})^2 T}/Z$,
which vanishes as $\epsilon \to 0$ since $Z > 0$ and the reference measure is finite.
This is a standard Gibbs concentration result: the measure concentrates exponentially
around the low-entropy manifold.
\end{proof}

The result follows the standard large-deviation phenomenon: the entropy weighting induces a
Gibbs measure that concentrates exponentially around low-entropy trajectories.  Since these
trajectories lie on a finite-dimensional manifold generated by the dynamical primitives,
biological computation occupies a sharply restricted subset of the full function space
despite its theoretical expressivity.

% ==========================================================
\section{Admissibility Constraints and the Revised Categorical Principle}
\label{app:admissibility}

\subsection{Motivation}

The biological prior of Appendix~\ref{app:cpg} constrains which trajectories occur.  A
complementary question concerns which \emph{transformations} are structurally valid: not
what a system does in practice, but what it is permissible for it to do while remaining
composable with other systems operating under the same constraints.  The answer takes the
form of a bounded region in the transformation space defined by the operator geometry
already in place, rather than a prescription over outputs.  A transformation is admissible
if and only if it remains within the region where phase alignment, entropy, and vorticity
are all bounded---precisely the conditions under which operators can be coherently composed.
This is the sense in which the framework yields something analogous to a categorical
constraint, but one that is grounded in geometry rather than derived from abstract ethical
reasoning.

\subsection{Admissible Transformation Set}

\begin{definition}[Admissible Operator and Admissible Set]
An operator $\mathcal{A}_x$ is \emph{admissible} at $x \in M$ if
\[
\theta(x) \le \theta_{\max}, \quad
S(x) \le S_{\max}, \quad
\xi(x) \le \xi_{\max}.
\]
The admissible set is
$\mathcal{A}_{\mathrm{adm}} = \{\mathcal{A}_x \in \mathcal{A} \mid \theta(x) \le
\theta_{\max},\; S(x) \le S_{\max},\; \xi(x) \le \xi_{\max}\}$.
\end{definition}

\subsection{Constraint Functional}

Define the global constraint energy
\[
\mathcal{C}[\gamma]
= \int_M \bigl[\alpha\,\theta(x)^2 + \beta\,S(x)^2 + \gamma\,\xi(x)^2\bigr]\,d\mu(x).
\]
A trajectory $\gamma$ is admissible if $\mathcal{C}[\gamma] \le C_{\max}$.

\subsection{The Revised Categorical Principle}

\begin{principle}[Admissibility Constraint]
A transformation is admissible if and only if it remains within bounded regions of phase
alignment, entropy, and vorticity such that it preserves composability with other admissible
transformations.
\end{principle}

This principle differs from classical categorical formulations in two ways.  First, it
operates at the level of transformations, not outputs: it does not prescribe what function
to compute, only which transformations are structurally valid.  Second, it implies both a
minimum and a maximum.  The lower bound---$\theta(x) \ge 0$, ensuring non-degenerate
transformation---prevents degenerate (null) operators that contribute no information to
the cascade.  The upper bounds---on $\theta$, $S$, and $\xi$---ensure stability and
composability.  An operator that violates the upper bounds produces non-composable
transformations, cascade breakdown, and unstable trajectories.

\subsection{Compositional Closure}

\begin{theorem}
If $\mathcal{A}_1, \mathcal{A}_2 \in \mathcal{A}_{\mathrm{adm}}$ and the constraint
functional $\mathcal{C}$ is subadditive over compositions, then
$\mathcal{A}_2 \circ \mathcal{A}_1 \in \mathcal{A}_{\mathrm{adm}}$.
\end{theorem}

\begin{proof}[Sketch]
Composition preserves the admissibility bounds when $\mathcal{C}$ satisfies the
subadditivity condition $\mathcal{C}[\mathcal{A}_2 \circ \mathcal{A}_1] \le
\mathcal{C}[\mathcal{A}_1] + \mathcal{C}[\mathcal{A}_2]$.
\end{proof}

\subsection{Relation to Biological Prior and Spectral Structure}

The admissibility set $\mathcal{T}_{\mathrm{bio}}$ from Appendix~\ref{app:cpg} can be
recovered as $\{\gamma \mid \mathcal{C}[\gamma] \le C_{\max}\}$, confirming that biological
systems operate within admissible bounds rather than maximising expressivity.  The
admissibility constraint also implies a spectral bound: $\sum_k \lambda_k a_k^2 \le E_{\max}$,
consistent with the spectral characterisation of $\mathcal{F}_{\mathrm{bio}}$ in
Appendix~\ref{app:spectral}.

\subsection{Alignment Without Identical Outputs}

The admissibility framing resolves a persistent ambiguity in alignment discussions.  Two
agents are aligned not because they produce identical outputs, but because their
transformations lie within the same admissible region and can be composed without breaking
coherence.  Alignment is a property of transformation geometry, not output agreement.  This
follows directly from the trajectory-equivalence account of propositional identity
(Appendix~\ref{app:equiv}): two agents that produce different surface responses may
nonetheless be $\Pi$-equivalent if their operator geometry is compatible under the same
perturbation sequence.

The Transformation Selection Principle from Section~\ref{sec:prior} can be restated in
this appendix's language.

\begin{principle}[Transformation Selection---Formal Statement]
\label{prin:selection_app}
Cognitive computation selects trajectories that minimise the constraint energy
\[
\mathcal{C}[\gamma]
= \int_M \bigl[\alpha\,\theta(x)^2 + \beta\,S(x)^2 + \gamma\,\xi(x)^2\bigr]\,d\mu(x)
\]
within the expressive envelope $\mathcal{F}_{\mathrm{all}}$, subject to task constraints.
\end{principle}

This formulation is not a prescription over outputs.  It is a constraint over the geometry
of admissible computation: cognitive systems do not compute any function they could in
principle compute, but the specific sub-class of functions reachable by trajectories that
remain within $\mathcal{A}_{\mathrm{adm}}$.

% ==========================================================
\section{Representation as Trajectory Equivalence}
\label{app:equiv}

\subsection{Equivalence Relation}

Let $\mathcal{X} = \{x \in M\}$ and let $\mathcal{P} = \{P_k\}$ be a sequence of
perturbations.  Define the perturbed cascade $x_{k+1} = \mathcal{A}_k(P_k(x_k))$.

\begin{definition}[Trajectory Equivalence]
Two states $x, y \in M$ are equivalent, written $x \sim y$, if for all admissible
perturbation sequences $\mathcal{P}$,
\[
\mathcal{A}^{(N)}_{\mathcal{P}}(x) = \mathcal{A}^{(N)}_{\mathcal{P}}(y)
\quad \text{for all } N.
\]
\end{definition}

\subsection{Quotient Space and Propositions}

The quotient $\mathcal{R} = \mathcal{X}/{\sim}$ represents propositions: an element
$[x] \in \mathcal{R}$ is a proposition.  A property $F\colon \mathcal{X} \to \mathbb{R}$
is representationally invariant if and only if $x \sim y$ implies $F(x) = F(y)$.

\subsection{Stability Under Composition}

\begin{theorem}
If $x \sim y$, then $\mathcal{A}^{(N)}(x) \sim \mathcal{A}^{(N)}(y)$ for any cascade.
\end{theorem}

\begin{proof}
Closure follows from associativity of operator composition.
\end{proof}

\subsection{Metric Structure}

Define the pseudometric
$d(x,y) = \sup_{\mathcal{P},N}\|\mathcal{A}^{(N)}_{\mathcal{P}}(x) -
\mathcal{A}^{(N)}_{\mathcal{P}}(y)\|$.
Then $x \sim y \iff d(x,y) = 0$.  Restricting to biologically admissible trajectories
yields ${\sim}_{\mathrm{bio}} \subseteq {\sim}$: biological systems identify propositions
only up to admissible transformations.

Propositions under this account are not static encodings.  They are invariants under
transformation cascades, and meaning is defined by stability under perturbation rather than
by any intrinsic property of the state.

% ==========================================================
\section{Empirical Observables and Measurement Protocol}
\label{app:measurement}

\subsection{Fixed Representation Manifold}

Let $M \subset \mathbb{R}^d$ be a pre-specified feature manifold with metric $g$, defined
prior to data collection (non-circular design).  All embeddings $x_t \in M$ are
constructed from this manifold; no parameters of $M$ are fitted to participant trajectories.

\subsection{Discrete Trajectory and Phase}

Given responses $\{x_t\}_{t=0}^T \subset M$, define $v_t = x_{t+1} - x_t$,
$\hat{v}_t = v_t/(\|v_t\|+\epsilon)$, and the discrete phase and amplitude
\[
\theta_t = \arccos\!\left(
\frac{v_t \cdot \nabla\Phi(x_t)}{\|v_t\|\,\|\nabla\Phi(x_t)\|+\epsilon}
\right),
\qquad
\rho_t = \|v_t\|,
\]
where $\Phi(x) = -d_M(x, x^\star)$ for the task target $x^\star$.

\subsection{Entropy Proxies}

The entropy field $S_t$ may be estimated by several proxies: the Shannon entropy of the
response choice distribution ($-\sum_i p_{t,i}\log p_{t,i}$), the surprisal of the
observed transition ($-\log p(x_{t+1}\mid x_t)$), or a monotone function of reported
confidence.  Any of these provide an operationalisation of $S$ that can be inserted into
the entropy-weighted operator.

\subsection{Vorticity, Aggregates, and Layer Segmentation}

Discrete vorticity is approximated as $\xi_t \approx \|\hat{v}_{t+1} - \hat{v}_t\|$.
Global aggregates are $\bar\theta = T^{-1}\sum_t \theta_t$ and $\Xi = \sum_t \xi_t$.
Partitioning indices into layers $\mathcal{I}_\ell$ (phonological, morphological,
semantic-pragmatic), layer-wise quantities are $\bar\theta^{(\ell)} =
|\mathcal{I}_\ell|^{-1}\sum_{t\in\mathcal{I}_\ell}\theta_t$ and
$\Xi^{(\ell)} = \sum_{t\in\mathcal{I}_\ell}\xi_t$.

A layer $\ell$ is \emph{convergent} if $\theta_{t+1} \le \theta_t + \delta$ for most
$t \in \mathcal{I}_\ell$ (small tolerance $\delta$).  At a layer boundary $b$, the
boundary vorticity $\xi_b = \|\hat{v}_{b+1} - \hat{v}_b\|$ measures the discontinuity
between cascade stages; large $\xi_b$ indicates a phase jump rather than smooth
composition.

\subsection{Human--Model Alignment Loss}

Let $\{x_t^{(h)}\}$ and $\{x_t^{(m)}\}$ be human and model trajectories embedded in
the same $M$.

\begin{definition}[Geometric Alignment Loss]
Let $\mathcal{A}^{(k)} = (\rho_k, \theta_k, S_k)$.  Define the layer-$k$ operator distance
\[
d_A^{(k)} = \sqrt{
\alpha_\rho(\rho_k^{(h)} - \rho_k^{(m)})^2
+ \alpha_\theta\, d_{S^1}(\theta_k^{(h)}, \theta_k^{(m)})^2
+ \alpha_S(S_k^{(h)} - S_k^{(m)})^2
},
\]
where $d_{S^1}(\theta, \theta') = \min(|\theta-\theta'|, 2\pi - |\theta-\theta'|)$ is the
angular geodesic distance on $S^1$, and $\alpha_\rho, \alpha_\theta, \alpha_S > 0$ are
weighting coefficients.  The alignment loss is $\mathcal{L}_A = \sum_k d_A^{(k)}$.
\end{definition}

Using $d_{S^1}$ corrects the periodicity of the phase variable and establishes
$\mathcal{L}_A$ as a proper geodesic distance on the product space
$\mathbb{R}_{>0} \times S^1 \times \mathbb{R}_{\ge 0}$.

Since $M$ and $\Phi$ are fixed prior to data collection, the estimates $\theta_t$ and
$\xi_t$ are independent of trajectory inference, ensuring that the test of phase alignment
is non-circular.  The primary experimental outputs are the phase trajectory $\{\theta_t\}$,
the vorticity profile $\{\xi_t\}$, the layer-wise convergence $\bar\theta^{(\ell)}$,
boundary discontinuities $\xi_b$, and the alignment loss $\mathcal{L}_A$.  Together these
define a fully operational pipeline mapping behavioural responses to the tuple
$(v_t, \theta_t, \rho_t, S_t, \xi_t)$ that constitutes the empirical amplitwist estimate.

% ==========================================================
\section{Non-Compositional Transformations and Failure Modes}
\label{app:failure}

\subsection{Compositional and Non-Compositional Transformations}

\begin{definition}[Compositional Transformation]
A transformation $f\colon M \to M$ is compositional if there exists a finite admissible
cascade $f \approx \mathcal{A}_N \circ \cdots \circ \mathcal{A}_1$ with each
$\mathcal{A}_k \in \mathcal{A}_{\mathrm{adm}}$.
\end{definition}

\begin{definition}[Non-Compositional Transformation]
$f$ is non-compositional if no finite admissible cascade approximates it with bounded
$\theta \le \theta_{\max}$, $\xi \le \xi_{\max}$, and $S \le S_{\max}$.
\end{definition}

\begin{proposition}
If $f$ is non-compositional, then any approximating trajectory $\gamma$ must contain a
time $t$ at which $\theta_t > \theta_{\max}$ or $\xi_t > \xi_{\max}$.
\end{proposition}

\subsection{Failure Signatures}

At a layer boundary $b$, the boundary jump is $\Delta v_b = \|\hat{v}_{b+1} - \hat{v}_b\|$.
Non-compositional tasks produce $\Delta v_b \gg 0$.  Biological systems minimising
$\mathcal{J}$ may respond by introducing spurious decompositions: a trajectory
$\gamma = \gamma_3 \circ \gamma_2 \circ \gamma_1$ where no true compositional structure
exists in $f$, reducing $\xi$ and $\theta$ locally at the cost of traversing an
intermediate state that is not part of the task.

\subsection{Three-Regime Prediction}

\begin{theorem}[Behavioural Regimes]
\label{thm:regimes}
Let tasks be partitioned into compositional and non-compositional families.  Then three
distinct patterns are predicted.  On compositional tasks, human trajectories show
$\theta_t \downarrow$ (improving phase alignment), $\xi_t \downarrow$ (stabilising
vorticity), and an ordered cascade in which lower layers stabilise before higher ones
engage.  On non-compositional tasks, human trajectories show $\theta_t \uparrow$ or
$\xi_t \uparrow$, or exhibit spurious intermediate stages.  Language models, on either
task family, exhibit low final output error alongside higher vorticity at layer boundaries
and disordered intermediate operator structure.
\end{theorem}

\subsection{Output vs Path Dissociation}

\begin{definition}[Endpoint Accuracy and Path Coherence Error]
Let $\epsilon_{\mathrm{out}} = \|f(x) - \hat{f}(x)\|$ and
$\epsilon_{\mathrm{path}} = \sum_t (\theta_t^2 + \xi_t^2)$.
\end{definition}

\begin{proposition}
It is possible that $\epsilon_{\mathrm{out}} \approx 0$ while $\epsilon_{\mathrm{path}} \gg 0$.
\end{proposition}

This distinguishes correct outputs from coherent transformations.  Non-compositional
transformations also excite high-frequency spectral modes ($a_k$ large for high
$\lambda_k$), simultaneously increasing entropy and vorticity.

A task is identified as non-compositional empirically if the observed data satisfy
$\Xi > \Xi_{\max}$ and $\bar\theta > \theta_{\max}$ despite low output error.  This
constitutes the falsifiable signature of the framework: the predicted asymmetry between
human failure modes (directional, prior-shaped, localised at decomposition boundaries)
and model failure modes (random, entropy-insensitive, uniformly distributed across layers)
is either present in the data or it is not.

% ==========================================================
\section{Gibbs Unification: Measure, Energy, and Admissibility}
\label{app:gibbs}

\subsection{Motivation}

The constraint functional $\mathcal{C}$, the entropy-weighted measure $\mu_{\mathrm{bio}}$,
and the admissibility region $\mathcal{A}_{\mathrm{adm}}$ are three descriptions of the
same object in different languages.  This appendix makes their identification explicit,
showing that admissible trajectories are precisely the typical set of a Gibbs distribution
whose energy functional is $\mathcal{C}$.

\subsection{Total Constraint Energy}

\begin{definition}[Total Constraint Energy]
Let $\gamma\colon[0,T]\to M$.  Define
\[
\mathcal{C}[\gamma]
= \int_0^T \Bigl(
\alpha\,\theta(\gamma(t))^2
+ \beta\,S(\gamma(t))
+ \delta\,\xi(\gamma(t))^2
\Bigr)\,dt.
\]
\end{definition}

\begin{definition}[Gibbs Measure on Trajectories]
\label{def:gibbs}
\[
d\mu_{\mathrm{bio}}(\gamma)
= \frac{1}{Z}\exp\!\bigl(-\lambda\,\mathcal{C}[\gamma]\bigr)\,d\mu_0(\gamma),
\]
where $\mu_0$ is a reference measure and $Z = \int \exp(-\lambda\mathcal{C})\,d\mu_0$
is the partition function.
\end{definition}

The entropy weighting of Appendix~\ref{app:entropy} and the constraint functional of
Appendix~\ref{app:admissibility} are now unified: $\mu_{\mathrm{bio}}$ is the Gibbs
measure induced by $\mathcal{C}$, and the exponential weight $w_k = e^{-\lambda S}$ is
the single-time-step restriction of the full path measure.

\subsection{Admissibility as a Typical Set}

\begin{definition}[Admissible Trajectory Set]
For $\epsilon > 0$, define
\[
\mathcal{T}_{\mathrm{adm}}^\epsilon
= \bigl\{\gamma \in \mathcal{T}_{\mathrm{all}} \mid
\mathcal{C}[\gamma] \le C_{\min} + \epsilon\bigr\},
\quad C_{\min} = \inf_\gamma \mathcal{C}[\gamma].
\]
\end{definition}

\begin{theorem}[Typical-Set Characterisation of Admissibility]
\label{thm:typical}
Let $\mu_{\mathrm{bio}}$ be as in Definition~\ref{def:gibbs}.  Then for any $\delta > 0$
there exists $\epsilon > 0$ such that
\[
\mu_{\mathrm{bio}}(\mathcal{T}_{\mathrm{adm}}^\epsilon) \ge 1 - \delta.
\]
\end{theorem}

\begin{proof}[Sketch]
Since $\mu_{\mathrm{bio}}(\gamma) \propto e^{-\lambda\mathcal{C}[\gamma]}$, trajectories
with energy exceeding $C_{\min} + \epsilon$ are suppressed by a factor of $e^{-\lambda\epsilon}$.
By standard large-deviation estimates, the measure concentrates on a thin shell around the
minimisers of $\mathcal{C}$, yielding the claim.
\end{proof}

\begin{corollary}[Variational Form of the Biological Prior]
Biological trajectories satisfy $\gamma \in \arg\min_{\tilde\gamma}\mathcal{C}[\tilde\gamma]$
almost surely under $\mu_{\mathrm{bio}}$.
\end{corollary}

In this formulation, the biological prior is not an external constraint imposed on
computation but the statistical consequence of a Gibbs distribution over trajectories.
Admissible trajectories are precisely those lying in the typical set of this measure.
Cognitive computation can therefore be understood as sampling from a low-temperature
Gibbs measure over transformation trajectories, where $\mathcal{C}$ defines the geometry
of admissible thought.

\subsection{Pushforward to Function Space}

\begin{proposition}[Concentration in Function Space]
\label{prop:fspace_conc}
Let $\Phi\colon\mathcal{T}_{\mathrm{all}}\to\mathcal{F}_{\mathrm{all}}$ map trajectories
to induced transformations.  Then the pushforward measure $\Phi_*\mu_{\mathrm{bio}}$
concentrates on a subset $\mathcal{F}_{\mathrm{bio}} \subset \mathcal{F}_{\mathrm{all}}$
of effective dimension bounded by the dimension of the minimiser set of $\mathcal{C}$.
\end{proposition}

Thus the restriction of biological computation to a low-dimensional function class is a
direct consequence of measure concentration in trajectory space.

% ==========================================================
\section{Variational Alignment and the Learning Objective}
\label{app:variational}

\subsection{Human and Model Gibbs Measures}

Let $\mathcal{C}_h$ and $\mathcal{C}_m$ denote the constraint functionals for human and
model systems respectively.

\begin{definition}[Induced Trajectory Measures]
\[
d\mu_h(\gamma) = \frac{1}{Z_h}e^{-\lambda\mathcal{C}_h[\gamma]}\,d\mu_0(\gamma),
\qquad
d\mu_m(\gamma) = \frac{1}{Z_m}e^{-\lambda\mathcal{C}_m[\gamma]}\,d\mu_0(\gamma).
\]
\end{definition}

Humans and models are treated as different Gibbs samplers over the same trajectory space,
distinguished by their effective constraint geometry.

\subsection{KL Divergence as the Canonical Alignment Quantity}

\begin{definition}[Geometric Alignment Divergence]
\[
D_{\mathrm{KL}}(\mu_h\,\|\,\mu_m)
= \int \log\!\left(\frac{d\mu_h}{d\mu_m}\right)\,d\mu_h.
\]
\end{definition}

Expanding using the Gibbs forms:
\[
D_{\mathrm{KL}}(\mu_h\,\|\,\mu_m)
= \lambda\,\mathbb{E}_{\mu_h}\bigl[\mathcal{C}_m - \mathcal{C}_h\bigr]
+ \log\frac{Z_m}{Z_h}.
\]
Misalignment is therefore the expected difference in constraint energy under human
trajectories.

\subsection{Alignment Loss as a KL Approximation}

\begin{proposition}[Alignment Loss as KL Surrogate]
\label{prop:klapprox}
Assume trajectories are observed at discrete steps $\{t_k\}$ and that $\mathcal{C}$
decomposes locally as a quadratic form in $(\rho, \theta, S)$.  Then
\begin{align*}
D_{\mathrm{KL}}(\mu_h\,\|\,\mu_m)
\;\approx\;
\mathbb{E}_{\mu_h}\!\Bigl[
\sum_k \bigl(
&\alpha_\rho(\rho_k^{(h)} - \rho_k^{(m)})^2
+ \alpha_\theta\,d_{S^1}(\theta_k^{(h)},\theta_k^{(m)})^2\\
&+ \alpha_S(S_k^{(h)} - S_k^{(m)})^2
\bigr)
\Bigr].
\end{align*}
\end{proposition}

\begin{proof}[Sketch]
Expanding $\mathcal{C}_m$ around $\mathcal{C}_h$ to second order yields a quadratic form
in local deviations of $(\rho, \theta, S)$.  Discretising the trajectory integral produces
the stated sum, with the angular geodesic $d_{S^1}$ arising from the periodicity of $\theta$.
\end{proof}

The alignment loss $\mathcal{L}_A$ defined in Appendix~\ref{app:measurement} is therefore
not an ad hoc metric but a Monte Carlo estimate of a KL divergence between trajectory
distributions.  Minimising $\mathcal{L}_A$ corresponds to matching the Gibbs measures
induced by human and model constraint functionals.  Alignment is achieved not when outputs
coincide, but when the two systems assign similar probability mass to the same regions of
trajectory space.

\subsection{Variational Alignment Objective}

\begin{proposition}[Variational Alignment Objective]
\label{prop:variational}
Let $\mu_h$ be fixed.  Minimising $D_{\mathrm{KL}}(\mu_h\,\|\,\mu_m)$ over model
parameters is equivalent to minimising
\[
\mathcal{J}_{\mathrm{var}}(\theta)
= \mathbb{E}_{\mu_h}\bigl[\mathcal{C}_m[\gamma]\bigr]
+ \frac{1}{\lambda}\log Z_m.
\]
\end{proposition}

This is a free-energy objective: expected energy plus log-partition.  Given sampled human
trajectories $\{\gamma^{(i)}\}$, the empirical approximation is
\[
\mathcal{J}_{\mathrm{var}} \;\approx\;
\frac{1}{N}\sum_{i=1}^N \mathcal{C}_m[\gamma^{(i)}]
+ \frac{1}{\lambda}\log Z_m,
\]
which, under the quadratic local approximation of Proposition~\ref{prop:klapprox},
reduces to the alignment loss $\mathcal{L}_A$ up to normalisation.

\subsection{Gradient of the Variational Objective}

\begin{proposition}[Gradient Form]
\label{prop:gradient}
\[
\nabla_\theta \mathcal{J}_{\mathrm{var}}
= \mathbb{E}_{\mu_h}[\nabla_\theta \mathcal{C}_m]
- \mathbb{E}_{\mu_m}[\nabla_\theta \mathcal{C}_m].
\]
\end{proposition}

The first term pushes the model to assign low constraint energy to human trajectories; the
second pushes energy up on the model's own typical trajectories.  This is structurally
identical to energy-based model training.

\subsection{Pathwise Alignment Criterion}

\begin{theorem}[Pathwise Alignment]
\label{thm:pathwise}
If $D_{\mathrm{KL}}(\mu_h\,\|\,\mu_m) = 0$, then for all admissible perturbation
sequences $\Pi$, the induced operator cascades satisfy
\[
\mathcal{A}^{(N)}_{\Pi,h}(x) = \mathcal{A}^{(N)}_{\Pi,m}(x)
\quad \text{almost surely under } \mu_h.
\]
\end{theorem}

Perfect alignment means identical transformation geometry, not merely identical outputs.
This reframes alignment from a problem of output matching to a problem of distributional
equivalence over transformation paths: two systems may agree on outputs while remaining
maximally misaligned in the geometry of their underlying trajectory measures.

Standard supervised learning minimises discrepancies in outputs via pointwise losses.  The
present objective minimises a divergence between distributions over trajectories.  Output
agreement is neither necessary nor sufficient for alignment; what is required is agreement
in the induced Gibbs measure over transformation paths.  Cognitive alignment is therefore
properly formulated as variational inference over trajectory space: the objective is to
learn a constraint functional whose induced Gibbs measure matches that of human cognition.
The amplitwist formalism provides the local parameterisation of transformations, the
constraint functional defines admissibility, and the alignment loss serves as a tractable
estimator of divergence between the resulting trajectory distributions.

% ==========================================================
\begin{thebibliography}{99}

\bibitem{needham1997visual}
T.~Needham,
\newblock \emph{Visual Complex Analysis}.
\newblock Oxford University Press, 1997.

\bibitem{ahlfors1979complex}
L.~V.~Ahlfors,
\newblock \emph{Complex Analysis}.
\newblock McGraw--Hill, 3rd edition, 1979.

\bibitem{conway1978functions}
J.~B.~Conway,
\newblock \emph{Functions of One Complex Variable}.
\newblock Springer, 2nd edition, 1978.

\bibitem{engel2000one}
K.-J.~Engel and R.~Nagel,
\newblock \emph{One-Parameter Semigroups for Linear Evolution Equations}.
\newblock Springer, 2000.

\bibitem{pazy1983semigroups}
A.~Pazy,
\newblock \emph{Semigroups of Linear Operators and Applications to Partial Differential Equations}.
\newblock Springer, 1983.

\bibitem{trotter1959product}
H.~F.~Trotter,
\newblock On the product of semi-groups of operators.
\newblock \emph{Proceedings of the American Mathematical Society},
  10(4):545--551, 1959.

\bibitem{kato1978trotter}
T.~Kato,
\newblock Trotter's product formula for an arbitrary pair of self-adjoint contraction semigroups.
\newblock \emph{Topics in Functional Analysis}, pp.~185--195, 1978.

\bibitem{evans2010pde}
L.~C.~Evans,
\newblock \emph{Partial Differential Equations}.
\newblock American Mathematical Society, 2nd edition, 2010.

\bibitem{chavel1984eigenvalues}
I.~Chavel,
\newblock \emph{Eigenvalues in Riemannian Geometry}.
\newblock Academic Press, 1984.

\bibitem{rosenberg1997laplacian}
S.~Rosenberg,
\newblock \emph{The Laplacian on a Riemannian Manifold}.
\newblock Cambridge University Press, 1997.

\bibitem{cheeger1970lower}
J.~Cheeger,
\newblock A lower bound for the smallest eigenvalue of the Laplacian.
\newblock \emph{Problems in Analysis}, pp.~195--199, 1970.

\bibitem{strogatz2018nonlinear}
S.~H.~Strogatz,
\newblock \emph{Nonlinear Dynamics and Chaos}.
\newblock CRC Press, 2nd edition, 2018.

\bibitem{izhikevich2007dynamical}
E.~M.~Izhikevich,
\newblock \emph{Dynamical Systems in Neuroscience}.
\newblock MIT Press, 2007.

\bibitem{marder2001central}
E.~Marder and R.~L.~Calabrese,
\newblock Principles of rhythmic motor pattern generation.
\newblock \emph{Physiological Reviews}, 76(3):687--717, 1996.

\bibitem{churchland2012neural}
M.~M.~Churchland et al.,
\newblock Neural population dynamics during reaching.
\newblock \emph{Nature}, 487:51--56, 2012.

\bibitem{gallego2017neural}
J.~A.~Gallego et al.,
\newblock Neural manifolds for the control of movement.
\newblock \emph{Neuron}, 94(5):978--984, 2017.

\bibitem{friston2010free}
K.~Friston,
\newblock The free-energy principle: a unified brain theory?
\newblock \emph{Nature Reviews Neuroscience}, 11:127--138, 2010.

\bibitem{tononi2008consciousness}
G.~Tononi,
\newblock Consciousness as integrated information.
\newblock \emph{Biological Bulletin}, 215(3):216--242, 2008.

\bibitem{carhart2014entropic}
R.~Carhart-Harris et al.,
\newblock The entropic brain: a theory of conscious states.
\newblock \emph{Frontiers in Human Neuroscience}, 8:20, 2014.

\bibitem{brandom1994making}
R.~Brandom,
\newblock \emph{Making It Explicit}.
\newblock Harvard University Press, 1994.

\bibitem{pearl2009causality}
J.~Pearl,
\newblock \emph{Causality: Models, Reasoning, and Inference}.
\newblock Cambridge University Press, 2nd edition, 2009.

\end{thebibliography}

\end{document}
